% Part: many-valued-logic
% Chapter: syntax-and-semantics
% Section: semantic-notions

\documentclass[../../../include/open-logic-section]{subfiles}

\begin{document}

\olfileid{mvl}{syn}{sem}

\olsection{Gagasan Semantis}

Upamakna logika mawa akeh nilai $\Log L$ diwenehake dening matriks.
Saiki kita bisa netepake gagasan semantis kang lumrah kanggo~$\Log L$.

\begin{defn} 
\begin{enumerate}
\item !!^a{formula}~$!A$ \emph{bisa dicukupi} yen ana
  $\pAssign{v}$ kang nyukupi, yaiku $\pSat{v}{!A}$; formula iku
  \emph{ora bisa dicukupi} yen ora ana $\pAssign{v}$ kang
  ndadekake $\pSat{v}{!A}$ kelakon;
\item !!^a{formula}~$!A$ iku \emph{tautologi} yen
  $\pSat{v}{!A}$ kanggo kabeh !!{valuation}s~$v$;
\item Yen $\Gamma$ himpunan !!{formula}s,
  $\Gamma \Entails !A$ (``$\Gamma$ ngakibatake $!A$'')
  yen lan mung yen $\pSat{v}{!A}$ kanggo saben
  !!{valuation}~$\pAssign{v}$ kang nyukupi $\pSat{v}{\Gamma}$.
\item Yen $\Gamma$ himpunan !!{formula}s, $\Gamma$
  \emph{bisa dicukupi} yen ana !!a{valuation}~$\pAssign{v}$
  kang nyukupi $\pSat{v}{\Gamma}$; yen ora, $\Gamma$
  \emph{ora bisa dicukupi}.
\end{enumerate} 
\end{defn}

Kanggo gagasan iki, ana sawetara fakta kang padha karo
kahanan ing logika klasik:

\begin{prop}
\ollabel{prop:semanticalfacts} 
\begin{enumerate} 
\item $!A$ iku tautologi yen lan mung yen
  $\emptyset \Entails !A$;
\item Yen $\Gamma$ bisa dicukupi, saben subhimpunan winates
  saka $\Gamma$ uga bisa dicukupi;
\item\ollabel{def:monotonicity}%
Monotonisitas: yen $\Gamma \subseteq \Delta$
  lan $\Gamma \Entails !A$, mula $\Delta \Entails !A$;
\item\ollabel{def:Cut}%
Transitivitas: yen $\Gamma \Entails !A$ lan
  $\Delta \cup \{ !A\} \Entails !B$, mula $\Gamma \cup \Delta \Entails
  !B$;
\end{enumerate}
\end{prop}

\begin{proof}
Latihan.
\end{proof}

\begin{prob}
Buktèkna \olref[mvl][syn][sem]{prop:semanticalfacts}.
\end{prob}

Ing logika klasik, kita bisa nggayutake akibat semantis karo kondisional.
Contone, \emph{modus ponens} iku sah: yen $\Gamma
\Entails !A$ lan $\Gamma \Entails !A \lif !B$, mula $\Gamma \Entails
!B$. Sesambungan wigati liyane antarane $\Entails$ lan $\lif$ ing
logika klasik yaiku teorema deduksi semantis: $\Gamma \Entails !A
\lif !B$ yen lan mung yen $\Gamma \cup \{!A\} \Entails !B$.
Asil-asil iki \emph{ora mesthi} lumaku ing logika mawa akeh nilai.
Lumaku-orane gumantung marang fungsi kabeneran~$\tf{\lif}$.

\end{document}
